\documentclass[aspectratio=169,11pt]{beamer}
\usetheme{Madrid}
\usecolortheme{dolphin}
\setbeamertemplate{navigation symbols}{}
\setbeamertemplate{caption}[numbered]

\usepackage[utf8]{inputenc}
\usepackage[T1]{fontenc}
\usepackage[spanish,es-noshorthands]{babel}
\usepackage{amsmath,amssymb,amsthm,mathtools}
\usepackage{tikz}
\usepackage{pgfplots}
\pgfplotsset{compat=1.18}
\usepackage{booktabs}
\usepackage{array}

\title[Prob. y Estadística]{Clase 0: Presentación de la materia}
\subtitle{Probabilidad y Estadística (5765)}
\institute[UNS]{Universidad Nacional del Sur \\ Departamento de Matemática}
\author{Probabilidad y Estadística (5765)}
\date{18 de agosto}

\begin{document}

\begin{frame}
\titlepage
\end{frame}

\begin{frame}{Contenidos}
\tableofcontents
\end{frame}

% ============================================================
\section{Bienvenida y objetivos}

\begin{frame}{Bienvenida}
\begin{center}
\Large ¡Bienvenidos a Probabilidad y Estadística!
\end{center}
\vspace{1em}
\begin{itemize}
\item Código 5765 — Área IV, Departamento de Matemática, UNS.
\item Correlativa cursada: Análisis Matemático II.
\item Carga horaria: 6 hs. teóricas + 6 hs. prácticas semanales.
\end{itemize}
\end{frame}

\begin{frame}{Objetivo de la materia}
\begin{block}{Del programa oficial}
Que el alumno comprenda y aplique los conceptos básicos de la probabilidad y de la estadística, conectándose con \emph{``el mundo de la aleatoriedad, de la inferencia y de la predicción''}, con fundamentos teóricos sólidos para comprender la utilidad y los límites de la estadística como herramienta en la resolución de problemas y en la toma de decisiones de la vida real.
\end{block}
\vspace{1em}
No es una materia de fórmulas para memorizar: es una materia para aprender a \textbf{pensar bajo incertidumbre}.
\end{frame}

\begin{frame}{Las 8 unidades del curso}
\begin{columns}[T]
\begin{column}{0.5\textwidth}
\begin{enumerate}
\item Espacios de probabilidad
\item Variables aleatorias
\item Vectores aleatorios
\item Función generadora de momentos
\end{enumerate}
\end{column}
\begin{column}{0.5\textwidth}
\begin{enumerate}
\setcounter{enumi}{4}
\item Muestras aleatorias y distribuciones
\item Estimación de parámetros
\item Prueba de hipótesis
\item Modelos lineales
\end{enumerate}
\end{column}
\end{columns}
\vspace{1em}
\centering 32 clases de teoría en total.
\end{frame}

\begin{frame}{Bibliografía}
\textbf{Básica}
\begin{itemize}
\item Meyer, P.: \emph{Probabilidad y Aplicaciones Estadísticas}. Addison Wesley Iberoamericana, 1992.
\item Rohatgi, V.: \emph{Statistical Inference}. John Wiley, 1984.
\item Walpole, Myers: \emph{Probabilidad y Estadística}. McGraw-Hill, 1992.
\end{itemize}
\textbf{De consulta}: Cánovos; Maronna (UNLP); Yohai; Bianco y Martínez (UBA) — disponibles en la web.
\vspace{0.5em}

\textbf{Software recomendado}: R (\texttt{www.r-project.org}), de uso público.
\end{frame}

% ============================================================
\section{Cronograma y modalidad}

\begin{frame}{Días y horarios}
\begin{itemize}
\item \textbf{Martes}: teoría — Clase 0 y Clase 1 el martes 18/8, y así sucesivamente.
\item \textbf{Jueves}: práctica (antes) y teoría.
\item Horario adicional: a confirmar.
\end{itemize}
\vspace{1em}
\begin{alertblock}{Inicio de clases}
Martes 18 de agosto.
\end{alertblock}
\begin{block}{Cronograma clase por clase}
El detalle día a día (fecha, martes/jueves y clase que corresponde, con chequeo de feriados) está en un documento aparte: \texttt{Cronograma.pdf}, en esta misma carpeta.
\end{block}
\end{frame}

\begin{frame}{Material de estudio}
Cada unidad del curso tiene, en el repositorio de la materia:
\begin{itemize}
\item \textbf{Teoría}: las clases en formato de diapositivas (beamer), una por tema.
\item \textbf{Un apunte de teoría extendido} por unidad: mismo contenido que las clases pero desarrollado en prosa, con más demostraciones y ejemplos, para quien quiera profundizar más allá de lo que da el horario de clase.
\item \textbf{Práctica}: guías de ejercicios resueltos paso a paso, de uso libre para estudiar.
\end{itemize}
\end{frame}

\begin{frame}{Asistencia}
\begin{itemize}
\item Las clases prácticas \textbf{no tienen asistencia obligatoria}: están para que cada uno practique a su ritmo.
\item Somos un grupo chico: si un día falta \textbf{todo el curso}, por favor avisen con anticipación.
\end{itemize}
\end{frame}

% ============================================================
\section{Evaluación}

\begin{frame}{Parciales}
\begin{itemize}
\item Se toman \textbf{2 parciales} durante el cuatrimestre.
\item Fechas: a confirmar con la asistente de cátedra (ver ubicación tentativa en el cronograma).
\item \textbf{La Unidad 8 (Regresión lineal y ANOVA) no entra en ningún parcial ni en el recuperatorio.}
\item Cada parcial se aprueba individualmente con \textbf{60 puntos} (sobre 100).
\item Hay \textbf{un (1) recuperatorio}, al final del cuatrimestre, para los parciales no aprobados.
\end{itemize}
\end{frame}

\begin{frame}{Condición final}
\begin{block}{Promoción directa}
Promedio de los 2 parciales \textbf{mayor a 80} $\Rightarrow$ promociona la materia sin final.
\end{block}
\begin{block}{Regular}
Con los 2 parciales aprobados (60 o más, incluyendo recuperatorio si hizo falta) pero promedio $\le 80$: queda regular y rinde final.
\end{block}
\end{frame}

\begin{frame}{Contacto y consultas}
\begin{itemize}
\item Profesor: \texttt{jmbavio@yahoo.com.ar}
\item Asistente de cátedra: (mail a confirmar)
\item Consultas: a demanda, coordinando por mail.
\end{itemize}
\end{frame}

% ============================================================
\section{Uso de Inteligencia Artificial}

\begin{frame}{Sobre el material del curso}
\begin{alertblock}{Transparencia sobre el material}
Gran parte de las clases de teoría, las prácticas resueltas y los apuntes extendidos de este curso fueron elaborados con asistencia de inteligencia artificial (IA), bajo la supervisión y revisión del docente responsable de la cátedra.
\end{alertblock}
\begin{itemize}
\item El contenido fue revisado, pero al ser una cantidad grande de material, \textbf{puede contener errores puntuales} (typos, algún cálculo, alguna definición mal formulada).
\item Si encontrás algo que no cierra: \textbf{avisá por mail}. Vamos a ir corrigiendo el material a lo largo del cuatrimestre.
\item El material es un complemento de estudio: no reemplaza la clase ni la bibliografía oficial.
\end{itemize}
\end{frame}

\begin{frame}{Uso de IA por parte de los alumnos}
\begin{itemize}
\item \textbf{Permitido} y \textbf{fomentado} para estudiar y practicar: usarla para entender un tema, revisar un ejercicio, generar problemas adicionales, etc.
\item \textbf{No permitido durante los parciales.}
\item Durante el cuatrimestre va a haber \textbf{hackatones y otras actividades} pensadas específicamente para aprender a usar la IA como herramienta de trabajo en probabilidad y estadística.
\end{itemize}
\vspace{1em}
\centering
\emph{La idea no es prohibir la herramienta, sino aprender a usarla bien.}
\end{frame}

% ============================================================
\section{Para arrancar}

\begin{frame}{Resumen}
\begin{itemize}
\item 8 unidades, 32 clases, martes y jueves desde el 18/8.
\item 2 parciales (60 aprueba c/u, recuperatorio único, promedio $>80$ promociona).
\item Prácticas libres, sin asistencia obligatoria.
\item Material generado con apoyo de IA y revisado por la cátedra — reporten erratas por mail.
\item IA permitida para estudiar, prohibida en parciales, protagonista en las hackatones.
\end{itemize}
\vspace{1em}
\begin{center}
\Large ¡Empezamos! Clase 1: Experimentos y espacio muestral.
\end{center}
\end{frame}

\end{document}
